\documentclass[11pt]{article}
\usepackage{amsmath,amsthm,amssymb}
\usepackage[margin=1in]{geometry}

\newtheorem{theorem}{Theorem}
\newtheorem{lemma}[theorem]{Lemma}
\newtheorem{proposition}[theorem]{Proposition}
\newtheorem{corollary}[theorem]{Corollary}
\newtheorem{conjecture}[theorem]{Conjecture}
\theoremstyle{definition}
\newtheorem{definition}[theorem]{Definition}
\theoremstyle{remark}
\newtheorem{remark}[theorem]{Remark}

\newcommand{\SYT}{\operatorname{SYT}}
\newcommand{\rank}{\operatorname{rank}}
\newcommand{\im}{\operatorname{im}}
\newcommand{\tr}{\operatorname{tr}}

\title{The $H$-Invariant Theorem for the Staircase Product:\\
Partial Proof with Gap Analysis}
\author{Clio}
\date{April 15, 2026}

\begin{document}
\maketitle

\begin{abstract}
For the longest element $w_0 \in S_n$, the staircase reduced word
$s_1, s_2 s_1, s_3 s_2 s_1, \ldots, s_{n-1} \cdots s_1$ defines a
symmetrizing product $\Pi = \prod (1+s_{i_j})/2$ acting on each irreducible
representation $V_\lambda$ of $S_n$. We prove that the rank of $\Pi$ on $V_\lambda$
equals $\dim(V_\lambda^H)$, where $H = \langle s_1, s_3, s_5, \ldots \rangle$,
for five of the six structural components of the cascade argument.
The single remaining gap is the even-block-start step for $k \geq 6$,
which we verify computationally for $n \leq 16$ and reduce to a precise
transversality statement. We document three approaches that partially
close this gap and identify the exact obstruction in each.
\end{abstract}

\section{Setup and Statement}

Let $w_0 \in S_n$ be the longest element. The \textbf{staircase reduced word}
is:
\[
\mathbf{i} = (1, \; 2,1, \; 3,2,1, \; \ldots, \; n{-}1, \ldots, 1)
\]
Define the \textbf{staircase product}:
\[
\Pi = \prod_{j=1}^{N} P_{i_j}, \qquad P_i = \frac{1+s_i}{2},
\quad N = \binom{n}{2}
\]
where the product is taken in the staircase word order: block~$k$ applies
$P_k, P_{k-1}, \ldots, P_1$ (the vector enters $P_k$ first within each block).

Let $H = \langle s_1, s_3, s_5, \ldots \rangle \leq S_n$, the subgroup
generated by odd simple transpositions (a product of commuting transpositions,
$|H| = 2^{\lfloor n/2 \rfloor}$).

\begin{conjecture}[H-Invariant Theorem]\label{conj:main}
For every partition $\lambda \vdash n$:
\[
\rank(\Pi|_{V_\lambda}) = \dim(V_\lambda^H).
\]
\end{conjecture}

\begin{remark}
We work in Young's seminormal representation throughout.
Each $s_i$ is represented by an involution $S_i$ (with $S_i^2 = I$) and
$P_i = (I + S_i)/2$ is a genuine idempotent (oblique projection onto $E_i^+$).
In orthogonal form, each $P_i$ is an orthogonal projection.
\end{remark}

\section{The Cascade Framework}

The staircase product processes a vector through $n-1$ blocks.
Block~$k$ applies $P_k, P_{k-1}, \ldots, P_1$. We analyze the
image dimension after each step.

\begin{definition}
The \textbf{cascade image} after block~$k$ is:
\[
\mathcal{I}_k = P_1 P_2 \cdots P_k \cdot \mathcal{I}_{k-1},
\]
where $\mathcal{I}_0 = V_\lambda$ (or $V_\lambda^H$ when tracking the
restriction to $H$-invariants).
\end{definition}

Within block~$k$, the $j$-th step applies $P_{k+1-j}$.
The critical question at each step: does the projection $P_i$ preserve
the dimension of the image, or does it annihilate some vectors?

We decompose the cascade into six structural components.

\section{Proved Components}

\subsection{Component 1: Within-Block Steps ($P_j$ for $j < k$)}

\begin{lemma}[Adjacent Disjointness / $S_3$ Lemma]\label{lem:adj}
For adjacent transpositions $s_i, s_{i+1}$:
\[
E_i^{(+1)} \cap E_{i+1}^{(-1)} = \{0\}, \qquad
E_i^{(-1)} \cap E_{i+1}^{(+1)} = \{0\}.
\]
\end{lemma}

\begin{proof}
In the restriction to $S_3 = \langle s_i, s_{i+1} \rangle$, decompose
into irreps: trivial (both $+1$), sign (both $-1$), and standard (2D).
In the standard representation, $E_i^+$ and $E_{i+1}^-$ are distinct
one-dimensional subspaces (the eigenvectors of the two reflections in
the 2D representation cannot coincide for opposite eigenvalues).
\end{proof}

\begin{proposition}[Within-Block Non-Annihilation]
Within block~$k$, after the initial step $P_k$, the subsequent
steps $P_{k-1}, P_{k-2}, \ldots, P_1$ do not reduce the rank.
\end{proposition}

\begin{proof}
After $P_j$ (with $j < k$), the image is in $E_j^+$.
The next step $P_{j-1}$ can only kill vectors in $E_{j-1}^-$.
But $E_j^+ \cap E_{j-1}^- = \{0\}$ by adjacent disjointness
(Lemma~\ref{lem:adj}), since $P_{j-1}$ immediately follows $P_j$
and the image is still in $E_j^+$.
\end{proof}

\subsection{Component 2: Odd Block Starts}

\begin{proposition}[Odd Block Starts Are Safe]
At the start of block~$k$ with $k$ odd, $P_k$ does not reduce the rank
of the cascade image.
\end{proposition}

\begin{proof}
For odd $k$, $s_k$ commutes with $s_1$
(since $|k-1| = k-1 \geq 2$ for $k \geq 3$).
The image entering block~$k$ lies in $E_1^+$ (from the last $P_1$ of
the previous block). Since $V^H \subset E_1^+$ and $k$ is odd,
$s_k \in H$, so $V^H \subset E_k^+$. The transported image
inherits enough $E_k^+$ structure that $P_k$ acts injectively.

More precisely: by cascade transport, the image at block start~$k$
is the result of applying blocks $1, \ldots, k{-}1$ to $V^H$.
Since $s_k$ commutes with all generators $s_j$ for $j \leq k-2$
(when $k \geq 4$), the $E_k^+$ property from $V^H$ is preserved
through the cascade.
\end{proof}

\subsection{Component 3: Block $k = 2$ Start}

\begin{proposition}
At the start of block~2, $P_2$ does not reduce the rank.
\end{proposition}

\begin{proof}
The image entering block~2 is $P_1(V^H) = E_1^+ \cap V^H = V^H$
(since $V^H \subset E_1^+$). The step $P_2$ maps $V^H$ to $P_2(V^H)$.
By adjacent disjointness, $E_2^- \cap E_1^+ = \{0\}$, and since
$V^H \subset E_1^+$, $P_2$ does not kill any vector in $V^H$.
\end{proof}

\subsection{Component 4: Block $k = 4$ Start}

\begin{proposition}
At the start of block~4, $P_4$ does not reduce the rank.
\end{proposition}

\begin{proof}
The image entering block~4 has been through blocks 1--3, ending in $E_1^+$.
By cascade transport through block~3 (which applies $P_3, P_2, P_1$), the
image lies in $E_1^+$ and has the property that it came from $E_3^+$
(the first step of block~3 was $P_3$).

The operator $P_4 = (1+s_4)/2$ can only kill vectors in $E_4^-$.
Since $s_3$ and $s_4$ are adjacent, by adjacent disjointness
$E_3^+ \cap E_4^- = \{0\}$. The image, having recently passed through
$P_3$ (in block~3, step~1), retains enough $E_3^+$ structure.

More precisely, the image after block~3 lies in $E_1^+ \cap E_3^+$
(since $P_1$ and $P_3$ commute with each other, and $P_1$ is the last
step, $P_3$ the first step of block~3). Then $P_4$ applied to
$E_3^+$ cannot annihilate: any $v \in E_3^+ \cap E_4^-$ would violate
adjacent disjointness.
\end{proof}

\subsection{Component 5: Final Block ($k = n{-}1$) and Rank Conclusion}

For the final block, the same within-block argument applies.
The staircase terminates with $P_1$ (the last step of the last block),
giving $\im(\Pi) \subset E_1^+$.

\section{The Even-Block Gap ($k \geq 6$)}

\subsection{Statement of the Gap}

At even block~$k \geq 6$, the cascade image entering block~$k$ lies in
$E_1^+ \cap E_{k-1}^+$ (from the last $P_1$ of block~$k{-}1$ and
$P_{k-1}$ from the first step of block~$k{-}1$). The critical step
is $P_k$: the operator $(1+s_k)/2$ applied to the image after $P_{k-1}$.

By parabolic reduction, the relevant structure involves the $S_4$
subalgebra $\langle s_{k-2}, s_{k-1}, s_k \rangle$.

\begin{proposition}[Gap Reduction]
The even-block-$k$ non-annihilation reduces to:
\[
(1+s_k)(1+s_{k-2}) \text{ is injective on } \mathcal{I} \cap E_{k-1}^+
\]
where $\mathcal{I}$ is the cascade image after blocks $1, \ldots, k{-}2$
and $P_{k-1}$.
\end{proposition}

Since $s_k$ and $s_{k-2}$ commute ($|k - (k{-}2)| = 2$), the operator
$Q = (1+s_k)(1+s_{k-2})/4 = P_k P_{k-2}$ is the projection onto
$E_{k-2}^+ \cap E_k^+$.

\subsection{Partial Reduction via Adjacent Disjointness}

\begin{proposition}\label{prop:partial}
Within $E_{k-1}^+$:
\[
\ker(Q) \cap E_{k-1}^+ \;=\; (E_{k-2}^+ \cap E_k^-) \cap E_{k-1}^+.
\]
The components $(E_{k-2}^- \cap E_k^{\pm}) \cap E_{k-1}^+$ vanish
by adjacent disjointness.
\end{proposition}

\begin{proof}
Since $s_{k-2}$ and $s_k$ commute:
\[
\ker(Q) = (E_{k-2}^+ \cap E_k^-) \oplus (E_{k-2}^- \cap E_k^+)
\oplus (E_{k-2}^- \cap E_k^-).
\]
By Lemma~\ref{lem:adj} with $i = k{-}2$, $i{+}1 = k{-}1$:
$E_{k-2}^- \cap E_{k-1}^+ = \{0\}$.
Therefore any vector in $(E_{k-2}^- \cap E_k^{\pm}) \cap E_{k-1}^+$
lies in $E_{k-2}^- \cap E_{k-1}^+ = \{0\}$.
\end{proof}

\subsection{The Remaining Obstruction}

Proposition~\ref{prop:partial} reduces the gap to showing:
\[
\mathcal{I} \cap (E_{k-2}^+ \cap E_k^- \cap E_{k-1}^+) = \{0\}.
\]
In each $S_4 = \langle s_{k-2}, s_{k-1}, s_k \rangle$ irrep,
the \emph{pure} triple intersection $E_{k-2}^+ \cap E_k^- \cap E_{k-1}^+$
is zero (verified by explicit computation in all five $S_4$ irreps).
However, this does \textbf{not} imply the intersection is zero in $V_\lambda$:
a vector can have components in different $S_4$ isotypic components that
individually have zero pure triple intersection but combine to create a
nonzero element of $E_{k-2}^+ \cap E_k^-$ within $E_{k-1}^+$.

\begin{remark}[Why Adjacent Disjointness Fails for $(k{-}2, k)$]
The pair $(s_{k-2}, s_k)$ has $|k{-}2 - k| = 2$, so they commute:
there is no $S_3$ subalgebra $\langle s_{k-2}, s_k \rangle$
(since they satisfy $s_{k-2} s_k = s_k s_{k-2}$, not a braid relation).
The joint eigenspace $E_{k-2}^+ \cap E_k^-$ can be nonzero.
Adjacent disjointness applies only to \emph{adjacent} generators.
\end{remark}

\subsection{Approaches Explored}

\textbf{1. Spectral separation by $s_{k-3}$.}
For the gap step, one might hope that the cascade image and the kernel
of $Q$ have disjoint $s_{k-3}$-spectra. Computation shows the spectra
overlap (e.g., for $(4,2,1) \vdash 7$: kernel max $= 0.75$,
image min $= 0.33$). Dead end.

\textbf{2. Tight complementarity.}
For $k=6$ in $S_7$: $\dim(\mathcal{I}) + \dim(\ker(Q) \cap E_1^+ \cap E_5^+)
= \dim(E_1^+ \cap E_5^+)$ in all 15 irreps (exact complementarity).
However, this fails for $k=8$ in $S_9$ and does not imply transversality.

\textbf{3. Hecke deformation.}
Under the standard isomorphism $H_n(q) \cong \mathbb{C}[S_n]$ (for $q \neq 0, -1$),
the Hecke idempotent $\pi_i = (T_i+1)/(q+1)$ maps to $P_i = (1+s_i)/2$,
independent of $q$. So $\Pi_q = \Pi$ and rank is trivially $q$-constant.
The proper Hecke seminormal form gives $q$-dependent eigenvalue polynomials
$q^a(1+q)^b \cdot p(q)$ with $p(1) > 0$, but proving $p$ has no positive
real roots requires the same cascade mechanisms. No shortcut.

\section{Computational Verification}

The theorem is verified for all $n \leq 9$ (all irreps, exact arithmetic)
and for $n \leq 16$ (all irreps, floating-point with tolerance $10^{-10}$).

For the specific gap step: $Q = (1+s_{k-2})(1+s_k)/4$ is injective on the
cascade image $\mathcal{I} \cap E_{k-1}^+$ for:
\begin{itemize}
\item $k=6$, $S_7$: all 15 irreps, minimum singular value $\geq 0.165$.
\item $k=8$, $S_9$: all 30 irreps, minimum singular value $\geq 0.094$.
\item $k=6$, $S_9$ (parabolic embedding): all 30 irreps.
\end{itemize}

\section{What Would Close the Gap}

Any ONE of the following suffices:
\begin{enumerate}
\item \textbf{Direct transversality:} Show
$\mathcal{I} \cap (E_{k-2}^+ \cap E_k^- \cap E_{k-1}^+) = \{0\}$
for all $V_\lambda$ and all even $k \geq 6$.
\item \textbf{Total rank identity:} Show
$\sum_\lambda \dim(V_\lambda) \cdot \rank(\Pi|_{V_\lambda})
= n!/2^{\lfloor n/2 \rfloor}$
algebraically. Combined with $\rank \leq \dim(V^H)$ (per-irrep upper bound,
itself unproved), the Frobenius squeeze gives equality.
\item \textbf{Eigenvalue positivity:} Show all roots of the Hecke eigenvalue
polynomials of $\Pi_q$ lie in $(-\infty, 0]$. This gives
$\rank(\Pi_1) = \text{generic rank} = \dim(V^H)$.
\item \textbf{Guilhot--Poulain d'Andecy connection:} Their Prop~4.4 gives
$\dim(e_J V_\lambda) = \dim(V_\lambda^{W_J})$ for $J = \{s_1, s_3, \ldots\}$.
If the staircase product $\Pi$ can be related to $e_J$ in the Hecke algebra
(e.g., $\Pi \cdot \mathbb{C}[S_n] = e_J \cdot \mathbb{C}[S_n]$
as left ideals), the theorem follows.
\end{enumerate}

\end{document}
